\section{Introduction} \label{sec:introduction}

According to a report released by the Organisation for Economic Co-operation and Development (OECD) in December 2023, Brazil remained below the average of OECD countries in PISA (\textit{Programme for International Student Assessment}), even though the country did not experience a statistically significant decline once the impact of the pandemic on Brazilian education is taken into account \parencite{oecd_pisa2022_brazil}. Learning deficits remain a long-standing bottleneck for the country's development and continue to motivate debates over the appropriate design of education policy.

In this context, policies that extend the school day have become increasingly prominent. In the state of São Paulo, the main example of this policy family is the Full-Time Education Program (\textit{Programa Ensino Integral}, PEI), established in 2012 \parencite{sao_paulo_pei_2012}. The program combines longer school hours with new management practices and the construction of each student's life project. Through staggered adoption, PEI expanded substantially over the last decade. Estimating the gains associated with policies such as PEI is important, but it also raises complex methodological challenges because of the way the policy was implemented.

The starting point of this paper is the empirical problem of estimating the gains associated with PEI on student performance in the São Paulo State School Performance Assessment System (SARESP). In a recent paper, \textcite{scorzafave2023} evaluated the program using the difference-in-differences (DiD) frontier estimator of \textcite{callawaysantanna} (CS21), overcoming the limitations of the traditional two-way fixed effects (TWFE) estimator, which, as shown by \textcite{goodmanbacon} and \textcite{chaisemartin2020}, can make ``forbidden comparisons'' and generate severely biased estimates under staggered adoption and heterogeneous treatment effects.

The structure of educational assessment data makes this problem even harder. SARESP exams do not follow the same student over time, forming repeated cross-sections (RCS). In addition, the empirical literature documents that converting a school to the PEI model is associated with changes in the composition of its student body \parencite{favaro_composicao_pei} and with spillover effects on nearby regular schools \parencite{santos2022_spillover_pei_escolas_regulares}. In short, the policy may change not only the treatment status of units, but also who appears in the sample after treatment. When this explicit compositional change occurs, most modern DiD estimators for RCS data, including \textcite{callawaysantanna}, rely on a stationarity assumption that may be strained, attributing to the program proficiency differences that could also reflect student selection \parencite{santannaxu}.

Against this background, this paper goes beyond replicating the empirical exercise in \textcite{scorzafave2023}. The central objective is to use PEI as an empirical laboratory to discuss which DiD estimator is most appropriate for this type of data, presenting a progression of estimators that address the structural problems with increasing rigor. I re-estimate the gains associated with PEI comparatively, moving through four approaches. First, I use the classical TWFE specification as a benchmark, while recognizing that it assumes homogeneous effects and fails structurally under staggered timing. Second, I apply the estimator of \textcite{callawaysantanna}, which solves the timing and temporal heterogeneity problem by estimating $ATT(g,t)$, replicating the baseline strategy of \textcite{scorzafave2023}. Third, because this method still imposes stationarity and does not correct for compositional changes in the cross-section, I use the doubly robust DiD estimator of \textcite{santannazhao} (DR-DiD), which combines outcome regression and weighting to handle covariates. Finally, I implement an aggregate approximation inspired by \textcite{santannaxu} (S\&X). This approximation does not fully reproduce the individual-level estimator proposed by the authors, because the available data do not allow individual enrollment records to be linked to individual SARESP scores; nevertheless, it allows me to assess whether incorporating observable compositional changes at the school-grade level changes the magnitude of the estimated effects.

The empirical strategy exploits the staggered adoption of the program across the state of São Paulo. I approximate the intuition of the \textcite{santannaxu} estimator to evaluate whether observable compositional changes alter the magnitude of the estimated effects. Pre-treatment dynamics are consistent with the parallel trends assumption in the preferred specifications, although they do not formally prove it. The estimates indicate that program adoption is associated with proficiency gains in standardized exams, both in Portuguese Language and Mathematics, for the two grades analyzed: 9th grade of elementary school and the 3rd grade of high school.

These results, however, do not fully separate learning gains, student selection, exam participation, and possible spillovers onto nearby regular schools. Because the data required to link enrollment records to exam microdata are unavailable, unlike the approach in \textcite{scorzafave2023}, the microdata are used only to construct aggregate covariates for the estimators. This limits the inferences that can be drawn from the results.

The contribution of this paper is to show why compositional changes matter for evaluations using RCS data, a common format in the economics of education and in standardized test databases, and to discuss through a transparent empirical adaptation whether the magnitude of the results changes when the observable composition of the student body is incorporated into estimation. The paper therefore contributes in two ways. First, in the economics of education, it re-evaluates PEI using estimators suited to staggered adoption and aggregate controls for student composition. Second, in applied causal inference, it makes explicit the distance between the ideal estimator, which would require perfectly linked microdata, and the feasible estimator in the available administrative setting.

The remainder of the paper is structured as follows. Section \ref{sec:literature} reviews the literature on full-time schooling, the structure of PEI, and the evolution of DiD estimators. Section \ref{sec:micromodel} presents the microeconomic model underlying the program's channels of action. Section \ref{sec:methodsanddata} describes the data and the theoretical intuition and assumptions behind the progression of econometric methods. Section \ref{sec:results} presents the estimation results, contrasts the estimators, and discusses the test for compositional changes. Finally, Section \ref{sec:conclusion} concludes.
